\begin{sourcescontribs}
  The main results of this chapter, namely \cref{thm-conds} and \cref{thm-uniquene-repr-minkowski}, were first published in \parencite{foundationsMDS}, then proved in \parencite{tversky1970dimensional} and finally refined and updated to the version presented in \parencite[Chapter 14.4.5, pp. 185-187]{foundations2}.
  This chapter is primarily based on the latter.\\
  The proofs presented in \cref{sec-conds} extend some arguments from \parencite[pp. 192-194]{foundations2} and include additional references.\\
  %

  The proofs presented in \cref{sec-unique-repr-thm} largely extend and complete the proofs provided in \parencite[pp. 195-199]{foundations2} and \parencite[p. 589-593]{tversky1970dimensional} \footnote{Until the end of this section, all references are with respect to these two sources, e.g., when writing "the proofs are briefly outlined"}. In both sources, the statements and proofs are almost identical.
  The proof of the main theorem, \cref{thm-uniquene-repr-minkowski}, can mostly be found in both sources.
  However, important details are missing in this proof. Most notably, the preconditions for the application of \cref{thm-power-mean-uniqueness} have neither been mentioned nor established.

  The proofs of the homogeneity theorem (\cref{sect-homogeneity-theorem}) and the convexity of $\metric$ balls (\cref{sec-ball-convexity}) are briefly outlined; however, some arguments are either skipped or only briefly mentioned without further explanation.
  Moreover, the statement of the homogeneity theorem is incorrect, as continuity is used but not assumed.
  \cref{lemma-part1} and \cref{lemma-part2} of the proof can be found in both references. The other parts of the proof are briefly outlined, so the majority of the work in this section involves the author completing the proofs.
  \cref{lemma-distance} was inspired by a post on math stackexchange, see \href{https://math.stackexchange.com/questions/48714/a-and-b-disjoint-a-compact-and-b-closed-implies-there-is-positive-dist}{https://math.stackexchange.com/questions/48714/a-and-b-disjoint-a-compact-and-b-closed-implies-there-is-positive-dist}.\\
  The definitions and convexity results in \cref{sec-ball-convexity} are based on Chapters 1.1 to 1.4 of \parencite[pp. 1-38]{lecturesConvGeo}. Most other arguments have been briefly outlined.

  The section on homogeneous quasi-arithmetic means, \cref{sect-HLP_proof}, is based on the work found in \parencite[pp. 66-69]{hardy1952inequalities} and \parencite[p. 271-273]{handbookMeans}. The proofs provided here can be found in these references.

\end{sourcescontribs}
